\section*{الفصل \ref{ch:g1:landscapes-around-us} --- المناظر من حولنا}

\begin{solution}{exo:g1:landscapes-around-us:1}
\omterm{def:g1:landscapes-around-us:landform}{الجبل}، \omterm{def:g1:landscapes-around-us:landform}{والتلة}، \omterm{def:g1:landscapes-around-us:landform}{والوادي}، \omterm{def:g1:landscapes-around-us:landform}{والسهل}.
\end{solution}

\begin{solution}{exo:g1:landscapes-around-us:2}
\omterm{def:g1:landscapes-around-us:landform}{الجبل}: فهو أعلى من \omterm{def:g1:landscapes-around-us:landform}{التلة} وأشد انحدارًا.
\end{solution}

\begin{solution}{exo:g1:landscapes-around-us:3}
طبيعي: \omterm{def:g1:landscapes-around-us:water}{النهر} والغابة والشاطئ. من صنع الإنسان: الطريق والجسر
والبيت.
\end{solution}

\begin{solution}{exo:g1:landscapes-around-us:4}
حيث تلتقي اليابسة بالبحر.
\end{solution}

\begin{solution}{exo:g1:landscapes-around-us:5}
مثلًا: طبيعي، \omterm{def:g1:landscapes-around-us:landform}{الجبال} أو الغابة على المنحدرات أو الجدول؛ ومن صنع
الإنسان، بيت أو كنيسة القرية.
\end{solution}

\begin{solution}{exo:g1:landscapes-around-us:6}
\omterm{def:g1:landscapes-around-us:landform}{الجبل}، وعلى قمته الثلج. \omterm{def:g1:landscapes-around-us:water}{والنهر} يجري في قاع \omterm{def:g1:landscapes-around-us:landform}{الوادي}، بين \omterm{def:g1:landscapes-around-us:landform}{الجبال}
\omterm{def:g1:landscapes-around-us:landform}{والتلة}.
\end{solution}

\begin{solution}{exo:g1:landscapes-around-us:7}
يجري \omterm{def:g1:landscapes-around-us:water}{النهر} من الأرض المرتفعة (\omterm{def:g1:landscapes-around-us:landform}{التلال} \omterm{def:g1:landscapes-around-us:landform}{والجبال}) نزولًا إلى \omterm{def:g1:landscapes-around-us:water}{بحيرة} أو إلى
البحر. فالماء ينحدر دائمًا، لذلك لا يستطيع أن يجري صاعدًا إلى قمة
\omterm{def:g1:landscapes-around-us:landform}{تلة}.
\end{solution}

\begin{solution}{exo:g1:landscapes-around-us:8}
الشجرة قريبة منا، في المستوى الأمامي؛ \omterm{def:g1:landscapes-around-us:landform}{والتلة} بعيدة، في المستوى الخلفي.
والأشياء البعيدة تبدو أصغر من الأشياء
القريبة.
\end{solution}

\begin{solution}{exo:g1:landscapes-around-us:9}
المستوى الأمامي: عشب طويل وبعض الشجيرات في أسفل الصورة.
المستوى الأوسط: \omterm{def:g1:landscapes-around-us:water}{النهر} يتعرج عبر \omterm{def:g1:landscapes-around-us:landform}{السهل} وعلى ضفتيه أشجار، وحقول، ومزرعة،
وطريق.
المستوى الخلفي: الأرض المستوية تمتد حتى الأفق، والسماء.
\end{solution}

\begin{solution}{exo:g1:landscapes-around-us:10}
لأن الناس صنعوه: أزالوا ما على الأرض، وحرثوها، وبذروا القمح. ولولا
الناس لكان هناك عشب أو شجيرات أو غابة، لا حقل قمح.
\end{solution}

\begin{solution}{exo:g1:landscapes-around-us:11}
لم يتغير: \omterm{def:g1:landscapes-around-us:landform}{التلة} \omterm{def:g1:landscapes-around-us:water}{والنهر}. وأضاف الناس: الحقول، وبيوتًا أكثر، والطريق،
والجسر (وقُطع معظم الغابة).
\end{solution}
